\subsection{Stable quotients in dimension five}

The complete component presentation in dimension six also gives a
quotient of the component group in dimension five. The next proposition
uses only the realization part of the classical obstruction theory.

\begin{proposition}\label{prop:five-dimensional-stable-surjectivity}
Let $M$ be a compact connected smooth five-manifold, possibly with
boundary. With the bottom and entire side boundary fixed and the
upper face free, positive suspension is surjective on components:
\[
 \sigma^+_*:\pi_0P_{\mathrm{Diff}}(M)
       \longrightarrow\pi_0P_{\mathrm{Diff}}(M\times I).
\]
Consequently every stable smooth concordance component is represented
on $M$. The corresponding assertions hold for compact PL
five-manifolds and their stipulated definable models over every real
closed field and in every fixed o-minimal expansion. For a finite
pointed homotopy model $K$ of the manifold, the
canonical stabilization map onto the complete presentation $E(K)$ is
a surjective homomorphism. Finite disjoint unions give finite products.

No hypothesis on $\pi_2M$ is required. Injectivity and a description
of the kernel are not asserted.
\end{proposition}

\begin{proof}
Write $C_5=\pi_0P_{\mathrm{Diff}}(M)$ and
$C_6=\pi_0P_{\mathrm{Diff}}(M\times I)$, rounding corners when
needed. As above, collar normalization and the fixed positive bending
give the homomorphism $s=\sigma^+_*:C_5\to C_6$. The first
obstruction is a surjective homomorphism
\[
 \Sigma_5:C_5\longrightarrow\operatorname{Wh}_2(\pi_1M).
\]
Its dimension-five surjectivity is exactly
\cite[Part~I, Chapter~VI, Theorem~2 and its surjectivity proof,
p.~213]{HatcherWagoner1973}. The proof constructs the elementary
handle-addition word and then cancels its identity intersection
matrix in dimension at least five. The first obstruction is
unchanged under positive suspension, by the remark on p.~130 of
the same source, so $\Sigma_6s=\Sigma_5$.

Only in dimension six do we use the complete corrected sequence
\eqref{eq:full-geometric-concordance-sequence}. Put
\[
 A=\operatorname{Wh}_1^+
       (\pi_1M;\mathbb Z/2\oplus\pi_2M),\qquad
 \mathcal R_\chi=\operatorname{im}\chi_{M\times I}.
\]
Thus $\ker\Sigma_6$ is the image of the injective map
$\overline\theta_6:A/\mathcal R_\chi\to C_6$. This includes the entire
Postnikov and framing correction. There is no dimension-five
exact sequence being assumed.

We claim that every element of $\overline\theta_6(A/\mathcal R_\chi)$ belongs
to $s(C_5)$. For this claim it suffices to realize any specified
finite coefficient sum before passing to the quotient $A/\mathcal R_\chi$.
Hatcher and Wagoner construct a marked one-parameter lens family
with any prescribed intersection coefficient in
$(\mathbb Z/2\oplus\pi_2M)[\pi_1M]$ whenever
$3\leq i\leq n-2$
\cite[Part~I, Chapter~VII, Proposition~5 and its complete
construction, pp.~227--228]{HatcherWagoner1973}.
Choose $n=5$ and $i=3$. This is the allowed endpoint $i=n-2$.

At that endpoint the normal bundle used in the local Whitney
construction has only one framed direction. The cited proof treats
this case separately: changing the inward pushing direction is
classified by $\pi_1S^1=\mathbb Z$, whose stabilization realizes
either value in $\pi_1O=\mathbb Z/2$. The choice of Whitney disk
realizes the specified $\pi_2M$ coefficient, and iteration gives
every finite sum with its fundamental-group labels. Thus the
endpoint has not been obtained by applying a higher-codimension
framing argument outside its range.

Positive suspension leaves the lens indices $3,4$ unchanged.
In the six-dimensional target these are in the stable marked-lens
range. The intersection curves, their maps to the loop space, and
their labels remain unchanged. The added positive coordinate is
a trivial line and preserves the stable framed bordism class. The stable convention
$\lambda^s_0=(-1)^3\lambda_0$ is compatible with this positive
suspension
\cite[Sections~4.1--4.2, pp.~8--9]{Igusa2021Dimension3}.
Since the source construction realizes every coefficient, the fixed
sign convention does not restrict its image. Forgetting the marked
lens data in the six-dimensional target gives exactly its geometric
$\theta_6$ map. Therefore each required coefficient class, and each
finite sum of them, has image $s(g)$ for a product of
five-dimensional lens concordances $g$. Passing to $A/\mathcal R_\chi$ can only
identify such images. This proves the claim without constructing a
five-dimensional homomorphism from $A/\mathcal R_\chi$ or classifying
five-dimensional lenses with equal invariants.

Now let $y\in C_6$. By surjectivity of $\Sigma_5$, choose
$x\in C_5$ with $\Sigma_5(x)=\Sigma_6(y)$. Then
$y\,s(x)^{-1}$ belongs to $\ker\Sigma_6$, hence by the claim
equals $s(g)$ for some $g\in C_5$. Thus $y=s(gx)$.
This proves surjectivity as a group statement and uses no splitting
of the corrected sequence.

The handle additions, disks, and lens modifications in the cited
realization occur in the interior. Collar normalization and positive
bending fix the entire side boundary, so the argument applies when
$M$ has boundary. The six-dimensional component theorem makes every
later stabilization an isomorphism. Hence the same source maps
surjectively onto the stable group.

A compact PL five-manifold admits a compatible smoothing, including
the boundary existence statement proved above. Let $N=M\times I$.
The smooth/PL component isomorphisms $\kappa_M$ and $\kappa_N$,
implemented through the common piecewise differentiable model,
respect composition and are compatible with forgetting to TOP.
The smooth and PL positive bends give the same map on relative
TOP components. Indeed, their marked parameter-disk straightenings
differ by an orientation-preserving homeomorphism preserving the
marked boundary arcs. Isotope its boundary-interval parameterizations
to the identity, extend over the disk, and apply the disk Alexander
isotopy to the remaining boundary-fixed map. Conjugation by this
isotopy preserves the pointwise bottom-and-side condition.
Only a TOP isotopy of the bends is used here.

The forgetful map
\[
 j_N:\pi_0P_{\mathrm{PL}}(N)\longrightarrow
          \pi_0P_{\mathrm{TOP}}(N)
\]
is injective because $\dim N=6$
\cite[Theorem~6.2, p.~41, and proof pp.~42--43]{BurgheleaLashof1974II}.
Consequently the equality
\[
 j_N\kappa_N\sigma^+_{\mathrm{Diff}}
   =j_N\sigma^+_{\mathrm{PL}}\kappa_M
\]
on relative components implies
$\kappa_N\sigma^+_{\mathrm{Diff}}
 =\sigma^+_{\mathrm{PL}}\kappa_M$.
Thus smooth surjectivity transfers to PL.
The relative finite-model PL/definable comparisons commute with
the fixed rational polyhedral suspension and with field extension,
so they transfer the same conclusion to each stipulated definable
model. The stable presentation identifies the target with $E(K)$
as a group under composition. Bottom fixation prevents permutation
of connected components, giving the finite-product assertion.
The empty case is immediate.
\end{proof}

The construction is a consequence of the cited realization theorems
and the corrected six-dimensional sequence. It does not resolve the
dimension-five injectivity question: equality of stabilized framed
coefficients does not by itself give a deformation of the original
lens families.
